\chapter{Introduction: what is verified, how to read this blueprint}
\label{chap:intro}

This blueprint documents a formal verification, carried out in Lean~4, of the erasure that the
\code{\#erase} command of this repository performs: the translation of a fully elaborated
\code{Lean.Expr} into a \lbox{} program --- the untyped calculus of Letouzey and of MetaCoq's
certified erasure --- which the \code{peregrine} middle-end then compiles to C, Wasm, OCaml,
CakeML, Rust or Elm. The subject of the verification is not an idealised model of that
translation but the shipping function itself, \code{Erasure.erase}
(\code{Erasure.lean:927}), together with the term walk \code{Erasure.visitExpr}
(\code{:590}) it is built from: the same code that the command elaborator calls and that writes
the \code{.ast} files the benchmark suite measures.

The development follows the architecture of Sozeau et al.'s certified erasure for Rocq
([S, Sec. 7.3--7.4]), which in turn machine-checks Letouzey's original design ([L]). That
architecture has four moving parts: a target calculus with a weak call-by-value big-step
semantics; a specification of \emph{erasability} and a non-deterministic erasure \emph{relation}
that the implemented erasure \emph{function} refines; a forward simulation for the relation; and
the collapse of the relation's non-determinism at first-order observations, which is what turns a
simulation into a statement about answers. All four are re-formalised here for Lean's own term
language and Lean's own kernel theory, the latter through \code{lean4lean}, whose \code{VEnv},
\TrExprS{}, \code{HasType} and \code{IsDefEq} are the only source of typing and
definitional-equality facts in this development and are never re-derived.

Two things distinguish the Lean setting from the Rocq one and force visible deviations,
catalogued in Section~\ref{sec:intro:correspondence}. First, Rocq writes constructors, case
analysis and fixpoints as \emph{term formers} (\code{tConstruct}, \code{tCase}, \code{tFix}),
while Lean writes all three as \code{Expr.const} heads: constructors are constants, pattern
matching goes through \code{casesOn} constants, and recursion is top-level. The erasure relation
\Erases{} is therefore a plain congruence over \code{Lean.Expr} with no case node and no fixpoint
node in its image at all, and the four Lean-specific compilation steps that produce those nodes
are factored out into a second relation \Lower{} between two \lbox{} terms. Second, the shipping
eraser reads the \emph{compiler-facing} body of a definition (the \code{\_unsafe\_rec} companion,
top-level recursive) rather than the kernel-facing recursor-based one, and it runs a preparation
pass before erasing; both are visible in the theorem as named hypotheses rather than as silent
assumptions. The main theorem, \code{shipping\_erase\_correct\_firstorder}
(\code{Capstone.lean:146}), is \textbf{conditional}: fifteen explicit binders, ten
of which stand as hypotheses at a general rung. This blueprint is written so that a reader can
see, node by node, which side of that line each statement falls on.

\section{The problem and the object of verification}
\label{sec:intro:problem}

The Peregrine project compiles programs written in proof assistants through one shared,
Rocq-verified middle-end whose input language is \lbox{}. A frontend's only job is to produce a
\lbox{} program; everything downstream is shared. This repository is the Lean frontend. Its user
interface is a single command, declared at \code{Erasure.lean:982} and elaborated
by \code{Erasure.eraseElab} (\code{:985}): \code{\#erase} takes a term, an optional configuration,
an optional output file and an optional \code{.mli} sidecar, elaborates the term, calls
\code{Erasure.erase} on it and serialises the resulting program.

The entry point is small enough to quote in words. \code{Erasure.erase} runs
\code{Erasure.visitExpr} on the result of \code{Erasure.prepare\_erasure}, starting from the empty
erasure state at the given configuration, and returns the untyped program consisting of the
accumulated global declarations and the erased main term, together with a list of inlining hints.
The two run hypotheses of the main theorem are exactly this decomposition --- one equation for the
preparation, one for the term walk --- and Chapter~\ref{chap:coldstart} owns the lemma
(\code{erase\_run\_ok}) that performs the split formally.

The preparation pass matters, because it is the pass that decides what the term walk actually
sees. \code{Erasure.prepare\_erasure} performs up to four rewrites: it redirects constants to
their \code{\_unsafe\_rec} companions (the function
\code{replaceUnsafeRecNames}), it applies Lean's own LCNF passes \code{macroInline} and
\code{inlineMatchers}, and --- only when the \code{csimp} flag is on --- it applies \code{@[csimp]}
replacements. The theorem's
syntactic conjuncts therefore speak about the \emph{prepared} term, while its observable conjunct
speaks about the term the user wrote, the two being connected by \code{prepare\_sound}
(Chapter~\ref{chap:coldstart}).

The configuration is a five-field record, \code{ErasureConfig}
(\code{Erasure.lean:64}), whose defaults are \code{extern := .preferAxiom}, \code{nat := .machine},
\code{csimp := true}, and \code{false} for both the
constructor-argument pruning and the typeclass-dispatch auto-inlining flags. Every theorem in this blueprint is stated under
\code{ConfigPinned} (\code{ErasureSpec.lean:43}), which demands
\code{csimp = false}, \code{extern = .preferLogical}, \code{nat = .peano} and the last two fields
false. \textbf{The pinned configuration differs from the default on three of the five fields}
(\code{csimp}, \code{extern}, \code{nat}): a default \code{\#erase} run is outside every theorem
stated here. The eight concrete instances of Chapter~\ref{chap:capstone} erase at
\code{Green.spikeConfig}, which is that pinned point.

Four things are deliberately outside the subject, and are listed once here and detailed in
Chapter~\ref{chap:trust}: the S-expression printer and the on-disk \code{.ast} grammar --- the
theorem's conclusion stops at the in-memory program; the \code{.ast.inlinings} and \code{.mli}
sidecars, and output size; the typed variant \lboxT{} that peregrine's Rust and Elm backends
require --- the conclusion names the untyped program; and machine \code{Nat}, \code{@[extern]}
axiomatisation and \code{@[csimp]}, all three excluded by \code{ConfigPinned}.

\begin{remark}
The design document of the \emph{subject} of this verification is [R], the internship report that
describes the shipping eraser, its preparation pass, its \code{Nat} and \code{Array} realizers and
its benchmark measurements. [R] contains no theorem and claims no correctness for the Lean side;
supplying one is what this development is for. Where [R] and the code disagree the code wins --- in
particular the tool it calls \code{lbox} is now \code{peregrine}.
\end{remark}

\section{The main theorem in one paragraph}
\label{sec:intro:theorem}

Suppose the shipping eraser is run on a source term $e$ from the empty erasure state under the
pinned configuration and returns the untyped program $(\Gamma, t)$; suppose the term it actually
walks --- the prepared term $pe$ --- lies in the supported fragment and translates into
\code{lean4lean}'s model of the Lean environment; and suppose the emitted program declares a body
for every constant it reaches and satisfies what peregrine's first pass reads of a program. Then
there are a \emph{specification environment} $\Gamma_{\mathrm{spec}}$ and an intermediate \lbox{}
term $t_0$ such that $pe$ erases to $t_0$ in the sense of the erasure relation \Erases{}; that
$\Gamma_{\mathrm{spec}}$ erases the source environment along everything $t_0$ reaches
(\ErasesEnv{}); that the lowering relation takes $t_0$ to the emitted term $t$ and
$\Gamma_{\mathrm{spec}}$ to the emitted environment $\Gamma$; and --- the observable --- for every
closed argument spine whose \lbox{} images carry erasures of their own, if the source evaluation
of $e$ applied to that spine yields a value $v$ whose type is a spine of an inductive that is
\emph{first order} in the model, then $v$ has a \textbf{unique} erasure, its lowering is
\textbf{box-free}, and the emitted program evaluates $t$ applied to the lowered arguments to
exactly that term under \lbox{}'s own weak call-by-value semantics.

That theorem is
\code{shipping\_erase\_correct\_firstorder} (\code{Capstone.lean:146}), stated in full as Theorem~\ref{thm:shipping_erase_correct_firstorder}
in Chapter~\ref{chap:capstone}. Three facts about it belong in the introduction.

\begin{itemize}
\item \textbf{It is conditional.} The theorem has fifteen explicit binders. Five of them ---
  the pinned configuration, the subject's translation, the fragment verdict, the body-less
  reference check and the output well-formedness check --- are discharged by checked terms at
  every one of the eight instances, and a sixth, the compiler-body hypothesis, at the first
  instance; the remaining ten are hypotheses. One of them, \code{hbridge} --- the two environment
  fields of \code{ErasureBridge} (\code{Capstone.lean:101}) --- is discharged \emph{nowhere, at any
  instance}: it is the single open proof obligation of the development.
\item \textbf{It is paired with eight concrete instances}, \code{Green.green\_G1} through
  \code{Green.green\_G8} (\code{LeanToLambdaBox/Green.lean}), each an instance of the theorem at a
  real \code{\#erase} run whose emitted program is committed in the file and whose conclusion ends
  in a literal Peano numeral, so that it cannot be satisfied by $\Box$ or by a stuck term. The
  environments they are stated at are universally quantified, so the ladder delivers
  \emph{conditional} non-vacuity; no computation makes it unconditional.
\item \textbf{What is deliberately not concluded}: MetaRocq's \code{expanded\_eprogram}, whose
  counterpart here is \code{LBExpandedFix}, and hence not \code{PeregrinePre}. The emitted body of
  every recursive declaration is a bare fixpoint node --- fifty of fifty measured --- so fixpoint
  eta-expansion is false on all five benchmark programs (shipping finding F-ETA). What the theorem
  does conclude is \code{LBWfPeregrine} (\code{Output.lean:227}), which the emitted
  programs do satisfy.
\end{itemize}

\section{The architecture of the proof}
\label{sec:intro:architecture}

\begin{enumerate}
\item \textbf{The target and its semantics} (Chapter~\ref{chap:lambdabox}). \code{LBTerm}
  (\code{Basic.lean:90}) is MetaRocq's \code{EAst.term} re-implemented for Lean,
  plus a free-variable constructor forced by the eraser's locally nameless representation and a
  primitive-value constructor, with branch binders as lists of binder names and no cofixpoints.
  Its semantics \WcbvEval{} (\code{Semantics/Eval.lean:79}) re-formalises MetaRocq's
  \code{EWcbvEval}, parameterised by a flag record; the correctness statement is made at the
  erasure flag point and carried to the entry flag point peregrine declares for its input. The
  certified evaluator with its soundness theorem \code{lbEval\_sound}
  (\code{Semantics/Compute.lean:497}) is what turns each instance's target-side evaluation into a
  kernel computation.
\item \textbf{The source side} (Chapter~\ref{chap:source}). Every typing, definitional-equality
  and translation fact comes from \code{lean4lean}, inherited verbatim. On top of it sits one
  source evaluation relation \SEval{}
  (\code{SourceEval.lean:193}) whose
  $\delta$ arm reads the \emph{compiler-body table} rather than the kernel body, with subject
  reduction \code{SEval.defeq}. The fragment the eraser is stated on is the predicate
  \code{Supported} (\code{Supported.lean:591}) with its sound Boolean checker
  \code{supportedB\_sound}.
\item \textbf{The erasure specification} (Chapter~\ref{chap:erases}). \Erasable{}
  (\code{Erasability.lean:55}) is [S]'s \code{isErasable}: a proof, or a type former (an arity
  \emph{up to definitional equality}). \Erases{} (\code{Erases.lean:211}) is [S, Fig. 18]
  transposed to \code{Lean.Expr}: eleven rules, ten congruences plus the one non-deterministic box
  rule, indexed by an environment, a level scope and a local context and nothing else --- no
  registry, no run state. \ErasesEnv{} (\code{ErasesEnv.lean:44}) is the dependency-selective
  environment erasure, [S]'s \code{erases\_deps}, in one clause with seven conditions.
\item \textbf{The forward simulation} (Chapter~\ref{chap:correctness}). \code{erases\_correct}
  (\code{Close.lean:29}) is [S]'s theorem of the same name: one induction on the
  source evaluation, with the $\iota$, $\delta$ and projection arms as separate step lemmas. It is
  stated at the \emph{emitted} environment and carries the \Lower{} factors, so there is exactly
  one simulation in this development and it is proved on the composite.
\item \textbf{Lowering} (Chapter~\ref{chap:lowering}). \Lower{} (\code{Lower.lean:336}), fourteen
  arms, relates two \lbox{} terms and is indexed by the specification environment alone. It
  carries the four Lean-specific compilation steps: constructor constants become constructor
  heads, saturated \code{casesOn} spines become case nodes, block members become fixpoint nodes,
  and the environment counterpart \LowerEnv{} prunes. The output boundary \code{LBWfPeregrine}
  (twelve clauses) is decided by a Boolean checker through \code{lbWfPeregrine\_of\_check}.
\item \textbf{The run model and the refinement} (Chapters~\ref{chap:eraser}, \ref{chap:run} and
  \ref{chap:refinement}). The module \code{ErasureRun.lean} (3620 lines) is the algebra of
  the eraser's monadic runs: what a successful run guarantees about the state, the name generator
  and the registries. On it, \code{visitExpr\_refines\_erasesLB} and its fixpoint companion
  (\code{VisitExprRefines.lean:190}, \code{:222}) say that a successful run of the shipping term
  walk lands in the composite \code{ErasesLB} --- this development's analogue of [S]'s
  \code{erases\_erase} and of [L, Lemma 11], factored because the shipping function is not a
  pruning function. The refinement is one eighteen-motive fixpoint induction whose member steps
  are hypotheses of the aggregator, supplied as proved terms by the capstone layer.
\item \textbf{The cold start and the bridge} (Chapter~\ref{chap:coldstart}).
  \code{erase\_run\_ok} splits the shipping entry point into its preparation and its term walk;
  \code{prepare\_sound} transports the source evaluation across the preparation; and
  \code{erasure\_bridge\_of\_run} (\code{Capstone.lean:63}) instantiates the bridge invariant at
  the empty entry state and applies the refinement with all eighteen steps supplied, concluding
  the erasure half. What it cannot supply is the specification environment's two \emph{environment}
  results; those are \code{ErasureBridge}'s two fields, and \code{bridgeEnv\_of\_regInv}
  (\code{Capstone.lean:122}) is the proved interface that would derive them from a registration
  invariant no theorem yet establishes for a run of the shipping eraser.
\item \textbf{The capstone and the instances} (Chapter~\ref{chap:capstone}). The composition, plus
  the eight instances. The capstone applies \code{erases\_correct} once at the spine, uses
  \code{firstorder\_erases\_core} for the answer's constructor-tree shape and its uniqueness, and
  \code{noBox\_lower\_of\_foSpine} to turn that shape into box-freedom of the lowered value.
\end{enumerate}

Everything above is proved \emph{from} the hypotheses collected in Chapter~\ref{chap:trust}: the
graph in this blueprint shows the proofs, and the assumption nodes of Chapter~\ref{chap:trust}
show where the proofs stop.

\section{Correspondence with the certified erasure of [S] and [L]}
\label{sec:intro:correspondence}

The two tables below map each notion of the reference development onto its counterpart here and
name the deviation. Line numbers are those of the tree this blueprint documents.

\begin{center}\small
\begin{tabular}{p{0.26\textwidth}p{0.30\textwidth}p{0.36\textwidth}}
\hline
Notion in [S] / [L] & This development & Deviation \\
\hline
\lbox{} syntax, [S, Fig. 16]; CIC-with-box, [L, Sec. 3.2] &
\code{LBTerm}, \code{Basic.lean:90} &
adds a free-variable constructor (the eraser is locally nameless) and a primitive one; branch
binders are a list of binder names; no cofixpoints, Lean having no coinduction \\
weak call-by-value evaluation, [S, Sec. 5.6 and 7.1]; weak reduction, [L, Def. 9] &
\WcbvEval{}, \code{Semantics/Eval.lean:79}, with a flag record, \code{Semantics/Flags.lean:36} &
big-step, not the closure of a one-step reduction; both constructor regimes are modelled and the
statement is pinned at applied form; the $\iota$ and projection rules drop MetaRocq's
parameter-count and saturation side conditions (Chapter~\ref{chap:lambdabox}) \\
box rule (1), a box applied to an argument reduces to a box, [S, Sec. 7.1]; [L, Def. 5] &
\code{WcbvEval.app\_box}, \code{Eval.lean:97} &
none: the argument is still evaluated and discarded, as in MetaRocq \\
box rule (2), singleton case on a box, [S, Sec. 7.1]; [L, Def. 8] &
\code{WcbvEval.iota\_sing}, \code{Eval.lean:171}, and \code{WcbvEval.proj\_prop}, \code{:200} &
gated on a flag \emph{and} on the inductive being marked propositional; the eraser never sets that
mark (measured false on 71 of 71 emitted inductives), so the rule is inert on emitted output ---
finding F-PROP \\
box rule (3), a fixpoint on a boxed guard, [S, Sec. 7.1]; [L, Def. 8] &
no separate rule: \code{WcbvEval.fix\_guarded}, \code{Eval.lean:210}, fires on the accumulated
argument count matching the principal argument index &
a boxed principal argument therefore unfolds like any other value; this realises the amendment
rather than omitting it \\
values, [S, Fig. 12] &
\code{Value}, \code{Semantics/Values.lean:85} &
applied-form constructor spines are values \\
\hline
\end{tabular}
\end{center}

\begin{center}\small
\begin{tabular}{p{0.26\textwidth}p{0.30\textwidth}p{0.36\textwidth}}
\hline
Notion in [S] / [L] & This development & Deviation \\
\hline
\code{isErasable}, [S, Sec. 7.2]; type scheme and proof, [L, Def. 1 and 3] &
\Erasable{}, \code{Erasability.lean:55} &
arity is taken up to definitional equality, which is closer to \code{Meta.isTypeFormerType} than a
syntactic test; Lean has no \code{Prop} into \code{Type} cumulativity, so [S]'s sort-quality
apparatus is unnecessary and is not built; proof irrelevance is definitional here and is used, not
re-proved \\
the erasure relation, [S, Fig. 18]; the invariant of [L, Def. 10] &
\Erases{}, \code{Erases.lean:211} --- eleven rules &
the variable rule splits in two (locally nameless); metadata and literal rules have no paper
counterpart; the constructor rule produces a \emph{bare head}, arguments arriving through the
application rule; \textbf{[S]'s case and fixpoint rules have no counterpart at all}, because
Lean's eliminators and recursion are constants, not term formers \\
the subsingleton side condition on the case rule, [S, Fig. 18 and Sec. 3.6.3] &
nothing in \Erases{}: the relation has no side condition anywhere at term level &
the content moves to the pass layer and to the fragment predicate; see F-PROP \\
environment erasure \code{erases\_deps}, [S, Sec. 7.4] &
\ErasesEnv{}, \code{ErasesEnv.lean:44} (one clause, seven conditions) &
the pointwise environment erasure [S] defines but does not use has deliberately no counterpart \\
the abstract environment, [S, Sec. 6.2] &
\code{ErasureSpec}, \code{ErasureSpec.lean:216} (eight fields), plus a reified source table and
its adequacy &
the queries are not squashed; adequacy of the reified slice is a named hypothesis mechanised
outside Lean rather than an interface property \\
\code{erases\_erase}, [S, Sec. 7.3]; [L, Lemma 11] &
\code{visitExpr\_refines\_erasesLB} and its fixpoint companion,
\code{VisitExprRefines.lean:190}, \code{:222}, landing in \code{ErasesLB} &
the shipping function is not a pruning function: it eta-expands to arity, splits \code{casesOn}
spines and builds literal towers, so the refinement is stated against the \emph{composite} of
\Erases{} and \Lower{}, not against \Erases{} alone \\
\code{erases\_correct}, [S, Sec. 7.4]; forward simulation, [L, Thm. 13] &
\code{erases\_correct}, \code{ErasesCorrect/Close.lean:29} &
stated at the emitted environment, with \Lower{} and \LowerEnv{} factors and an extra upstream
hypothesis bundle; the source evaluation is hypothesised, Lean having no progress or
standardization counterpart here \\
\code{firstorder\_ind}, [S, Sec. 7.4]; data type and logic-free type, [L, Def. 14 and 6] &
\code{FirstOrderInd}, \code{FirstOrderInd.lean:86} &
MetaRocq's shipped predicate is cited as origin and not transcribed: its sort conjunct makes it
false on the natural numbers (reported upstream); the monomorphism and index-freedom fields here
are declared scope restrictions booked to neither paper \\
determinism of erasure at a first-order type, [S, Sec. 7.4] &
\code{firstorder\_erases\_deterministic}, \code{FirstOrderInd.lean:592}; in the capstone,
\code{firstorder\_erases\_core} with \code{noBox\_lower\_of\_foSpine} &
the collapse is stated for the relation against itself (uniqueness of the erasure), not against an
erasure function, since no erasure function is defined here \\
\code{axiom\_free}, [S, Sec. 5.6] &
\code{NoBodylessRefs}, \code{Output.lean:862} &
read on the emitted environment, decidable, and a premise of the capstone only; the simulation
needs no analogue, a body-less constant having no source value \\
first-order erasure correctness, [S, Sec. 7.4]; [L, Thm. 15] &
\code{shipping\_erase\_correct\_firstorder}, \code{Capstone.lean:146} &
stated in $\Downarrow$ form off a hypothesised source evaluation, not off reduction to an
irreducible term; the observable is quantified over a closed first-order \emph{argument spine}
([S]'s separate-compilation corollary folded in); concludes \code{LBWfPeregrine} and deliberately
not \code{LBExpandedFix} \\
\code{optimize} and its correctness, [S, Sec. 7.4] &
\code{LBOptimize\_correct}, \code{Optimize.lean:785} &
off the shipping path: stated at block flags, and F-PROP makes the pass inert on this frontend's
output --- a worked example, not a pipeline pass \\
the four Lean-specific compilation steps &
\Lower{}, \code{Lower.lean:336}, and \LowerEnv{}, \code{ErasesEnv.lean:227} &
\textbf{no counterpart in [S] or [L]}; anchored instead by this repository's own
\code{doc/rules-Lower.md} against MetaRocq's iota-reduction and fixpoint substitution \\
--- &
\code{Supported}, \code{Supported.lean:591}, and \code{ConfigPinned},
\code{ErasureSpec.lean:43} &
\textbf{no counterpart}: a decidable fragment predicate over the term \emph{and its dependency
closure}, and the configuration restrictions, both theorem-visible \\
--- &
\SEval{}, \code{SourceEval.lean:193} &
\textbf{no counterpart}: [S] reuses PCUIC's own evaluation, whereas \code{lean4lean} has none, so
one is defined here; its $\delta$ arm reads the compiler-body table \\
\hline
\end{tabular}
\end{center}

Five of these deviations are load-bearing and are taken up in the chapters shown.

\begin{itemize}
\item \textbf{Applied-form constructors} (Chapter~\ref{chap:lambdabox}). peregrine's first pass
  requires applied form; the erasure flag point pins the non-block regime and the constructor rule
  of \Erases{} emits a bare head. [L] has no such distinction at all --- constructors are context
  variables there --- and [S] does not discuss it; it is a CertiCoq-era concern. Measured: all 982
  emitted constructor nodes carry an empty argument list.
\item \textbf{No case rule and no fixpoint rule in \Erases{}} (Chapter~\ref{chap:erases}). This is
  the structural consequence of Lean writing constructors, inductive type names and constants all
  as \code{Expr.const}. The module's own docstring puts it exactly: ``Rocq writes a constructor
  \code{tConstruct}, an inductive type name \code{tInd} and everything else \code{tConst}; Lean
  writes all three \code{Expr.const}, so the relation says which.''
\item \textbf{No emitted fixpoint is applied} (Chapter~\ref{chap:lowering}). The eraser emits a
  fixpoint node as a whole constant body, so MetaRocq's \code{expanded\_eprogram} is false on all
  five benchmark programs; that is finding F-ETA, and it is exactly the gap between the concluded
  \code{LBWfPeregrine} and the defined but unconcluded \code{PeregrinePre}.
\item \textbf{The first-order fragment} (Chapters~\ref{chap:correctness} and~\ref{chap:trust}).
  Both [S] and [L] restrict the \emph{observation}, not the program; so does this development. But
  first-orderness of the natural numbers is a hypothesis at all eight instances, and the
  predicate's only proved model inhabitant in this tree is at an empty inductive.
\item \textbf{Machine \code{Nat}, \code{@[csimp]} and \code{@[extern]}}
  (Chapter~\ref{chap:trust}). All three are shipping features excluded by \code{ConfigPinned};
  [R, Sec. 4.2] describes what they do and why they exist, and no theorem in this blueprint covers
  a run with any of them enabled.
\end{itemize}

\section{How to read this blueprint}
\label{sec:intro:legend}

\paragraph{Node kinds.} Five standard theorem-like environments are used --- \emph{definition}
(types, relations, executable functions), \emph{lemma}, \emph{proposition}, \emph{theorem},
\emph{corollary} --- and one non-standard one added for this blueprint, \emph{assumption}: a
statement the development assumes rather than proves. Assumption nodes appear in
Chapter~\ref{chap:trust} only; other chapters point at them. Remarks, like this paragraph's
surroundings, are prose and are not graph nodes.

\paragraph{Graph colours.} They are \code{leanblueprint}'s, not this blueprint's invention. A
node's \emph{border} shows the status of its statement: green when the statement is formalised
(it names Lean declarations that exist and elaborate), blue when every prerequisite is formalised
but the statement itself is not, orange when marked not-ready. A node's \emph{fill} shows the
status of its proof: green when the Lean proof is complete in this repository, blue when it is
only ready to be attempted, dark green when the node and every ancestor are proved.

\paragraph{What that means here, precisely.} The repository contains no \code{sorry} of its own,
so every result node in Chapters~\ref{chap:lambdabox}--\ref{chap:capstone} is green in both
senses: its statement elaborates and its proof is complete. \textbf{Green therefore means ``proved
in Lean from the hypotheses of its own statement''. It does not mean ``unconditionally true'', and
for the capstone it emphatically does not mean that the erasure is correct.} The hypotheses live
in Chapter~\ref{chap:trust} as assumption nodes, which carry their Lean names on the statement
when a Lean declaration states them and carry \emph{no proof environment at all}, so the graph
renders them as stated-but-unproved and every path into the capstone visibly passes through one.

\paragraph{The four hypothesis bundles.} \code{ErasureSpec} (eight fields: specifications of
Lean's \code{Meta} and \code{Core} primitives and of the model connection); \code{EraserAsks}
(five fields, obligations about this repository's own code, two of which, on the kernel's
inductive head and on distinctness of block keys, are refuted in general by
findings F-DEPTH and F-UNSAFEREC and hold only on all measured inputs); \code{UpstreamAsks} (four
fields, \code{lean4lean} lemmas that a pin bump would supply); and \code{ErasureBridge} (two
fields, the open problem). Alongside them stand the theorem's inline binders, named in
Chapter~\ref{chap:trust} one by one.

\paragraph{The \code{lean4lean} trust boundary.} Every typing, definitional-equality and
term-translation fact on the source side is \code{lean4lean}'s. The pin is the head of the
\code{trproj} branch of a fork, recorded in \code{lakefile.toml}; it carries sixteen \code{sorry}
sites, enumerated in the committed fixture \code{test/lean4lean-sorries.expected} and re-measured
by \code{scripts/lean4lean-sorries.sh}. Results that reach them inherit \code{sorryAx}. This is
the development's one inherited trust boundary, in the sense in which [S] inherits a normalization
axiom, and it is an assumption node with no Lean counterpart in Chapter~\ref{chap:trust}.

\paragraph{Measurement, not narration.} Axiom footprints are measured, not asserted:
\code{test/ledger.expected} holds the literal \code{\#print axioms} output of forty invocations in
\code{test/Ledger.lean}, and CI diffs it. The capstone and each of the eight instances carry one
and the same set of 33 names, including \code{sorryAx} and twenty-nine non-standard reflection
axioms of Lean's executable kernel --- the price of discharging the relevance oracle through
\code{lean4lean}'s executable checker rather than assuming its soundness.
\code{visitExpr\_refines\_erasesLB} and its fixpoint companion, \code{bridgeEnv\_of\_regInv},
\code{LBOptimize\_correct} and \code{lbEval\_sound} are measured free of \code{sorryAx}.
Chapter~\ref{chap:trust} carries the tables and explains what \code{\#print axioms} does and does
not see: it measures a proved theorem and never a hypothesis, so moving an obligation into a
binder moves no row.

\paragraph{One further reading rule.} Where a statement in this blueprint is more specific than
its prose --- an extra flag, a closedness side condition, a pinned configuration --- the chapter
says so in one sentence after the informal statement. Where the development deviates from [S] or
[L], the deviation is named in the text, never left to the reader to notice.

\section{Reading order}
\label{sec:intro:order}

\begin{itemize}
\item \emph{A reader who knows [S, Sec. 7] and wants the delta.}
  Section~\ref{sec:intro:correspondence}, then Chapter~\ref{chap:erases} (the relation),
  Chapter~\ref{chap:lowering} (the part [S] does not have), Chapter~\ref{chap:capstone} (the
  theorem) and Chapter~\ref{chap:trust} (what it assumes).
\item \emph{A reader who wants to know whether to trust the output of \code{\#erase}.}
  Section~\ref{sec:intro:theorem}, then Chapter~\ref{chap:trust} first and
  Chapter~\ref{chap:capstone} second --- in that order, because the theorem's content is bounded
  by its hypotheses.
\item \emph{A reader following the proof.} Chapters~\ref{chap:lambdabox}
  and~\ref{chap:source} (the two languages), \ref{chap:erases} and~\ref{chap:correctness}
  (specification and simulation), \ref{chap:lowering} (the emitted form),
  \ref{chap:eraser}--\ref{chap:refinement} (the shipping function, its run algebra, the
  refinement), \ref{chap:coldstart} (the cold start), \ref{chap:capstone} (composition and
  instances) and \ref{chap:trust} (the boundary).
\end{itemize}

The repository carries tracked anchors this blueprint does not duplicate and does not replace:
\code{doc/rules-Erases.md} (the rule-by-rule transport table against [S, Fig. 18]),
\code{doc/rules-Lower.md}, \code{doc/trust.md} and the generated \code{doc/coverage.md}; and the
check commands \code{lake exe green-check --all}, \code{lake exe reify --check},
\code{lake exe coverage} and \code{bash scripts/ledger.sh}.

\section{References}
\label{sec:intro:references}

\textbf{[S]} Matthieu Sozeau, Yannick Forster, Meven Lennon-Bertrand, Jakob Botsch Nielsen,
Nicolas Tabareau and Th\'eo Winterhalter, \emph{Correct and Complete Type Checking and Certified
Erasure for Coq, in Coq}, Journal of the ACM 72(1), Article 8, January 2025. Section 7 is the
certified erasure; the sections used here besides it are 5.6 (weak call-by-value standardization,
\code{axiom\_free}, values), 3.6.3 (allowed eliminations) and 6.2 (the abstract environment). A
local copy is \code{references/MetaCoq\_paper.pdf}. Citations in this blueprint are by section, so
[S, Sec. 7.3].

\textbf{[L]} Pierre Letouzey, \emph{A New Extraction for Coq}, in Types for Proofs and Programs
(TYPES 2002), Lecture Notes in Computer Science, Springer, pages 200--219. This is the paper that
introduces the calculus, there called CIC with a box and now universally called \lbox{}, the
erasure function, the box reduction and the correctness statement for closed terms of a logic-free
data type. A local copy is \code{references/A\_New\_Extraction\_for\_Coq.pdf}. Two cautions this
blueprint observes: Letouzey never writes ``\lbox{}'', and the forward simulation --- one source
step matched by one target step or by box steps --- is his Theorem 13, not Theorem 12, which runs
the other way and is what transports termination.

\textbf{[R]} Simon Dima, \emph{Compiling Lean programs with Rocq's extraction pipeline}, MPRI
internship report, supervised by Yannick Forster (Inria Paris, Cambium), 1 October 2025. A local
copy is \code{references/lean\_extraction\_report.pdf}. It is the design document of the subject
of this verification, not of its specification.

MetaRocq's own sources (\code{EAst}, \code{EWcbvEval}, \code{ExAst},
\code{EOptimizePropDiscr}, \code{PCUICFirstorder}) are cited by file wherever a rule-by-rule
comparison is made, since [S] does not display the \lbox{} evaluation relation at all; the on-disk
interchange format is peregrine-tool's \code{doc/format.md} and \code{theories/PAst.v}.

\paragraph{Deviations from the paper and caveats.} This chapter has claimed four things that a
reader of [S] should carry into the rest of the blueprint. The erasure relation here has no case
rule and no fixpoint rule, and the compilation steps that build those nodes live in a second
relation, \Lower{}, that has no counterpart in either paper. The correctness statement is made
against the emitted environment and off a \emph{hypothesised} source evaluation: nothing here
proves that a Lean term evaluates, only that if it does, the emitted program agrees with it. The
target semantics is faithful to MetaRocq rule for rule only at the non-block flag point, and the
$\iota$ and projection rules drop MetaRocq's parameter-count and saturation side conditions.
Finally, one binder of the capstone, \code{hnb}, is bound and never used in its proof, and the
block hypothesis \code{hblk}, while forwarded to \code{erasure\_bridge\_of\_run}, is spent by no
theorem in the closure --- so two of the theorem's premises are stronger than its proof requires.
Chapter~\ref{chap:trust} records these, and every other finding, with its measurement.
